\section{Lista de verificación de ingeniería: convertir un prototipo de investigación en un sistema de producción}

\subsection{Diez preguntas que deben responderse antes del lanzamiento}
En la etapa de prototipo de investigación, el sistema suele buscar mostrar casos de éxito lo antes posible; pero al pasar a producción, es necesario anteponer la robustez, el costo y la capacidad de auditoría. La siguiente lista puede usarse como plantilla de revisión previa al lanzamiento.

\begin{table}[H]
\centering
\small
\begin{tabular}{p{0.8cm}p{3.9cm}p{5.8cm}}
\toprule
\textbf{\#} & \textbf{Pregunta} & \textbf{Ejemplo de criterio de aprobación} \\
\midrule
1 & ¿Los límites de la tarea son claros? & Las restricciones de entrada, las acciones prohibidas y las condiciones de terminación están todas definidas en la documentación \\
2 & ¿El estado puede reproducirse? & Cada observation/action/result de cada paso puede reproducirse y auditarse \\
3 & ¿El presupuesto de búsqueda es controlable? & Se configuran un número máximo de pasos, una ramificación máxima y una estrategia de reversión por tiempo de espera agotado \\
4 & ¿El Judge está calibrado? & La precisión alcanza el objetivo en el conjunto de control fuera de línea, y existe un mecanismo de revisión periódica \\
5 & ¿La falla puede recuperarse automáticamente? & Admite una recuperación de tres etapas: detección de anomalías, reversión y replanificación \\
6 & ¿El disparador de intervención humana está definido con claridad? & Existe una estrategia de escalamiento obligatorio para escenarios de pago, privacidad y baja confianza \\
7 & ¿El retorno de datos forma un ciclo cerrado? & Tanto las trayectorias exitosas como las fallidas pueden retornar y participar en la siguiente ronda de entrenamiento \\
8 & ¿El costo puede explicarse? & Cada tipo de tarea cuenta con estadísticas de costo de tokens, latencia y llamadas externas \\
9 & ¿La evaluación se actualiza de manera continua? & Tanto el benchmark como el tráfico real se monitorean de manera continua \\
10 & ¿La seguridad y el cumplimiento están cubiertos? & Las operaciones sensibles cuentan con privilegio mínimo y registros de auditoría \\
\bottomrule
\end{tabular}
\caption{Lista mínima de ingeniería antes de la conversión en producto de un Web Agent}
\end{table}

\subsection{Los tres tipos de costos de ingeniería que más se subestiman}
\begin{enumerate}
    \item \textbf{Costo de gestión de estado}: un prototipo sin estado es rápido, pero la complejidad aumenta abruptamente al incorporar reversión y estado con múltiples ramificaciones;
    \item \textbf{Costo de evaluación}: construir un conjunto de control de alta calidad y calibrar el Judge de manera continua toma más tiempo que entrenar una demo;
    \item \textbf{Costo de operación}: la estructura de las páginas web cambia con frecuencia, por lo que las reglas de análisis y el mapeo de acciones requieren mantenimiento continuo.
\end{enumerate}

\begin{knowledgebox}{Por qué la ``operación'' se convierte en el costo principal}
Muchos equipos, en la etapa de PoC, solo observan el desempeño del modelo e ignoran la falla de las reglas provocada por los cambios en las páginas web. Lo que realmente consume mano de obra después del lanzamiento suele ser problemas de operación continua, como cambios de selector, rediseños del flujo de la página y actualizaciones de los mecanismos anti-automatización.
\end{knowledgebox}

\subsection{De impulsado por métricas a impulsado por valor}
El ponente mencionó ``economically valuable tasks'', lo que nos recuerda que no basta con mirar la puntuación del benchmark, sino que hay que incorporar el valor de la tarea al objetivo de optimización. Una tarea con puntuación alta pero valor bajo no debería ocupar un presupuesto de inferencia considerable; por el contrario, una tarea de alto valor debería permitir un mayor costo de búsqueda y una colaboración humana más intensa.

\begin{importantbox}{Recomendación de programación escalonada por valor}
\begin{itemize}
    \item Tareas de bajo valor: estrategia ligera + presupuesto bajo + falla rápida;
    \item Tareas de valor medio: búsqueda moderada + reversión automática;
    \item Tareas de alto valor: búsqueda intensiva + verificación múltiple + intervención humana.
\end{itemize}
\end{importantbox}

\subsection{Resumen de la sección}
La dificultad central de la conversión en producto de un Agent no está en ``si puede ejecutar una demo'', sino en ``si puede operar de manera continua, estable, controlable y auditable''. La lista de ingeniería y la programación escalonada por valor son las palancas clave para superar este umbral.

